% !TEX root = main.tex
% joint_draft.tex --- DRAFT (not included in main.tex).
% Analysis of joint volume + centroid formulations and of the role of the
% conservation identities (eq:partition-feasibility) as hard constraints.
\documentclass[main.tex]{subfiles}

% Local notation (move to alias.tex if adopted)
\providecommand{\momset}{\mathcal{M}}            % moment set of all (fractional) partitions
\providecommand{\mom}{\mathbf{m}}                % moment vector of a cell, \mom_i = (\momx_i, \cellvol_i)
\providecommand{\tmom}{\overline{\mathbf{m}}}    % prescribed moments, \tmom_i = (\pvol_i \tcentroid_i, \pvol_i)
\providecommand{\momx}{\mathbf{s}}               % first moment of a cell, \momx_i = \cellvol_i \centroid_i
\providecommand{\gen}{\mathbf{z}}                % homogeneous generator, \gen_i = (\dual_i, \theta_i/2)
\providecommand{\support}{\sigma_{\momset}}      % support function of \momset
\providecommand{\convset}{\mathcal{C}}           % centroid slice of \momset at fixed volumes
\providecommand{\regdual}{\jointdual_{\alpha\beta}} % proximal joint dual
\providecommand{\tgt}{\mathbf{t}}                % shifted centroid target of the proximal problem
\title{Joint constraints draft}  % standalone build only (ignored when included)

\begin{document}

\section{Joint volume and centroid constraints, revisited}
\label{sec:joint-revisited}

This section re-examines the jointly constrained problem of \S\ref{sec:joint-constraints}.
We first show that $\jointdual$ is, up to a constant, the negated support function of a convex set of cell moments, which explains both its degeneracy and its frequent unboundedness.
We then ask whether the conservation identities \eqref{eq:partition-feasibility}, imposed as hard constraints, can restore feasibility, and show that they cannot.
Finally, we propose a \emph{proximal} joint dual that keeps both constraint types, is strictly concave, and reduces to known constructions in its limits.

\subsection{The moment set}
\label{sec:moment-set}

A \emph{fractional partition} of $\domain$ is a collection of functions $\rho_i \colon \domain \to [0,1]$ with $\sum_i \rho_i \equiv 1$; partitions are the special case of indicator functions $\rho_i = \basis_i$.
We collect the zeroth and first moments of each $\rho_i$ into $\smash{\mom_i = \bigl( \momx_i, \cellvol_i \bigr) = \bigl( \int \rho_i \, \xx \, d\xx, \int \rho_i \, d\xx \bigr) \in \R^{d+1}}$, so that $\momx_i = \cellvol_i \centroid_i$.
The \emph{moment set} $\momset \subset \R^{n(d+1)}$ is the set of all such moment vectors $\mom = \{\mom_i\}$.
It is convex and compact, as the linear image of the convex compact set of fractional partitions.
Summing over $i$ shows that $\momset$ lies in the affine subspace defined by \eqref{eq:partition-feasibility}, which has dimension $(n-1)(d+1)$, and one can check that $\momset$ has full dimension within it.
Its support function is
\begin{equation}
  \support(\gen)
  = \max_{\mom \in \momset} \sum_i \gen_i\tr \mom_i
  = \int_\domain \max_i \bigl( \dual_i\tr \xx + \tfrac{1}{2}\theta_i \bigr) \, d\xx ,
  \qquad \gen_i = \bigl( \dual_i, \tfrac{1}{2}\theta_i \bigr) ,
  \label{eq:support}
\end{equation}
because the maximum over fractional partitions of a pointwise-linear objective is attained by assigning each $\xx$ to the largest affine function.
The maximizing partition is the power diagram with generators $\dual_i$ and weights $\theta_i + \norm{\dual_i}^2$, and it is unique (up to null sets) whenever the generators are distinct.

\paragraph{$\jointdual$ is a support function.}
Expanding the shifted power distance of \eqref{eq:complete-square} as $\norm{\xx}^2 - 2\bigl(\dual_i\tr\xx + \tfrac12\theta_i\bigr)$ with $\theta_i = \weight_i - \norm{\qq_i}^2 - 2\,\shift_i\tr\tcentroid_i$, and substituting $\weight_i$ back into the term $\sum_i \weight_i \pvol_i$ of \eqref{eq:joint-dual}, gives
\begin{equation}
  \jointdual(\shift, \weight)
  = \kappa - 2 \Bigl( \support(\gen) - \sum_i \gen_i\tr \tmom_i \Bigr) ,
  \qquad
  \kappa = \sum_i \pvol_i \norm{\qq_i - \tcentroid_i}^2 + \int_\domain \norm{\xx}^2 d\xx - \sum_i \pvol_i \norm{\tcentroid_i}^2 ,
  \label{eq:joint-support}
\end{equation}
where $\tmom_i = (\pvol_i \tcentroid_i, \pvol_i)$ are the prescribed moments and the map $(\shift, \weight) \mapsto \gen$ is affine and invertible.
The constant $\kappa$ is exactly the transport cost $\sum_i \int_{\cell_i} \norm{\xx - \qq_i}^2 d\xx$ of \emph{any} partition with volumes $\pvol_i$ and centroids $\tcentroid_i$, which confirms that the primal objective carries no information once both constraints are imposed.
Since $\support$ is convex and positively $1$-homogeneous, \eqref{eq:joint-support} immediately yields three facts.

\begin{proposition}[Structure of $\jointdual$]
\label{prop:joint-structure}
Assume the data $\tmom$ satisfy \eqref{eq:partition-feasibility}.
\begin{enumerate}
  \item $\jointdual$ is concave, and $\jointdual - \kappa$ is positively $1$-homogeneous in $\gen$.
  \item Its Hessian, wherever it exists, annihilates the $(d\!+\!2)$-dimensional space spanned by the $d\!+\!1$ modes of \eqref{eq:joint-nullspace} and by the radial direction $\delta\gen = \gen$, \ie $\delta\dual_i = \dual_i$ and $\delta\theta_i = \theta_i$, a uniform dilation of the shifted sites about the origin (dilations about other points differ by a translation mode).
  \item $\sup \jointdual = \kappa$ if $\tmom \in \momset$, attained at the degenerate point $\gen = 0$ (all generators coincide), and $\sup \jointdual = +\infty$ otherwise.
  The supremum is attained with distinct generators if and only if $\tmom$ is an exposed point of $\momset$, \ie if and only if some power diagram realizes the prescribed volumes and centroids.
\end{enumerate}
\end{proposition}

\begin{proof}[Sketch]
(1) is inherited from $\support$.
(2) The $d\!+\!1$ modes of \eqref{eq:joint-nullspace} change $\support(\gen) - \gen\tr\tmom$ by $\translation\tr\bigl( \int \xx - \sum_i \pvol_i \tcentroid_i \bigr) + c \bigl( \domainVol - \sum_i \pvol_i \bigr) = 0$; the radial mode is in the kernel by Euler's identity for $1$-homogeneous functions.
(3) $\support(\gen) - \gen\tr\tmom \ge 0$ for all $\gen$ iff $\tmom \in \momset$ (separation), and by homogeneity the infimum is then $0$; equality at $\gen \neq 0$ means $\gen$ exposes $\tmom$, and the exposed point is the moment vector of the power diagram of $\gen$.
\end{proof}

Proposition~\ref{prop:joint-structure} corrects \S\ref{sec:joint-constraints} in two places: the null-space is $(d\!+\!2)$-dimensional (the dilation of Lemma~2.1 in \citet{Bourne2025Inverting} survives on a fixed domain), and the maximizer is not unique up to these modes, since $\gen = 0$ is always a maximizer.
After eliminating the weights, $\max_\weight \jointdual$ is (up to a factor and a constant) the helper function $g$ of \citet[Eq.~(4.1)]{Bourne2025Inverting} and the function $H(\cdot\,;B)$ of \citet[Eq.~(40)]{Bourne2026ConvexOrder}, and items (2)--(3) are their Theorem~4.2(v),(vii) and Theorem~3.6, respectively.

\paragraph{Realizable data has codimension one.}
Power diagrams are parametrized by $n(d+1)$ generator coordinates modulo the $(d\!+\!2)$-dimensional redundancy of item~(2), so the realizable moments form a set of dimension at most $n(d+1) - d - 2$.
This is one less than the dimension $(n-1)(d+1)$ of the affine subspace \eqref{eq:partition-feasibility}: realizable data lie on the relative boundary of $\momset$.
For instance, in $\domain = [0,1]$ with two cells of volume $\tfrac12$, the slice of $\momset$ is the segment $\tcentroid_1 \in [\tfrac14, \tfrac34]$, $\tcentroid_2 = 1 - \tcentroid_1$, and only its two endpoints are realizable \citep[Example~4.7]{Bourne2025Inverting}.
The same count with the volumes fixed gives $nd - d - 1$ for the realizable centroids, matching the relative boundary of the convex set $\convset(\pvol) = \{ \centroid : (\pvol_i \centroid_i, \pvol_i)_i \in \momset \}$, whose exposed points are exactly the realizable centroids \citep[Thm.~3.6]{Bourne2026ConvexOrder}.

\subsection{What the conservation identities can and cannot do}
\label{sec:feasibility-hard}

The identities \eqref{eq:partition-feasibility} are $d\!+\!1$ \emph{affine} conditions that cut out the affine hull of $\momset$.
Realizability, in contrast, is the \emph{non-affine} condition of lying on the relative boundary of $\momset$, so no set of affine constraints can enforce it.
Three variants of ``imposing \eqref{eq:partition-feasibility} as hard constraints'' make this precise.
\begin{itemize}
  \item \emph{On fixed data.} The identities only test that $\tmom$ lies in the affine hull of $\momset$; if they fail, $\jointdual$ is unbounded along the modes \eqref{eq:joint-nullspace}, and if they hold, the data are still generically non-realizable.
  \item \emph{On slack variables.} Relaxing the constraints to $\cellvol_i = \pvol_i + \varsigma_i$ and $\momx_i = \pvol_i\tcentroid_i + \mathbf{r}_i$, and imposing $\sum_i \varsigma_i = 0$ and $\sum_i \mathbf{r}_i = 0$, is vacuous: every fractional partition satisfies them when the data do. Their multipliers are a global translation and a global constant, \ie exactly the gauge modes \eqref{eq:joint-nullspace}, and the radial mode of Proposition~\ref{prop:joint-structure} is untouched.
  \item \emph{On free unknowns.} When some targets become unknowns (\eg free volumes for prescribed centroids), the identities define the polytope $\{ \cellvol \ge 0 : \sum_i \cellvol_i = \domainVol,\ \sum_i \cellvol_i \tcentroid_i = \domainVol\centroid_\domain \}$ of the max-entropy remark in \S\ref{sec:remarks}. It is an \emph{outer} approximation of the convex set of admissible volumes $\{\cellvol : ((\cellvol_i\tcentroid_i, \cellvol_i))_i \in \momset\}$; it is useful for initialization and as a cheap infeasibility certificate, but it is implied by the formulation below and does not select realizable data.
\end{itemize}
A well-posed joint formulation must therefore replace the exact targets by a \emph{projection} onto $\momset$, which lands on the relative boundary whenever the target lies outside.

\subsection{Proximal joint dual}
\label{sec:proximal-joint}

For parameters $\alpha, \beta > 0$ we define the \emph{proximal joint dual}
\begin{equation}
  \regdual(\shift, \weight)
  = \jointdual(\shift, \weight)
  - \frac{\alpha}{2} \sum_i \pvol_i \norm{\shift_i}^2
  - \frac{\beta}{2} \sum_i \pvol_i \, \weight_i^2 .
  \label{eq:proximal-dual}
\end{equation}

\begin{proposition}[Primal problem]
\label{prop:proximal-primal}
Maximizing $\regdual$ is dual to the convex problem over fractional partitions
\begin{equation}
  \min_{\mom \in \momset} \
  \sum_i \int \rho_i \norm{\xx - \qq_i}^2 d\xx
  + \frac{2}{\alpha} \sum_i \frac{\norm{\momx_i - \cellvol_i \tcentroid_i}^2}{\pvol_i}
  + \frac{1}{2\beta} \sum_i \frac{(\cellvol_i - \pvol_i)^2}{\pvol_i} ,
  \label{eq:proximal-primal}
\end{equation}
with no duality gap, and the optimal multipliers are
\begin{equation}
  \shift_i = \frac{2}{\alpha} \frac{\cellvol_i}{\pvol_i} \bigl( \tcentroid_i - \centroid_i \bigr) ,
  \qquad
  \weight_i = \frac{1}{\beta} \frac{\pvol_i - \cellvol_i}{\pvol_i} .
  \label{eq:proximal-multipliers}
\end{equation}
\end{proposition}

\begin{proof}[Sketch]
The Lagrangian \eqref{eq:joint-lagrangian}, extended to fractional partitions, is affine in $\mom$ and in $(\shift, \weight)$; subtracting the strongly concave regularizer of \eqref{eq:proximal-dual} keeps it concave in the multipliers, and $\momset$ is convex and compact, so Sion's minimax theorem applies.
Maximizing over $\shift_i$ the term $-2\shift_i\tr(\momx_i - \cellvol_i\tcentroid_i) - \frac{\alpha}{2}\pvol_i\norm{\shift_i}^2$ gives the centroid penalty, and maximizing over $\weight_i$ the term $-\weight_i(\cellvol_i - \pvol_i) - \frac{\beta}{2}\pvol_i\weight_i^2$ gives the volume penalty; the maximizers are \eqref{eq:proximal-multipliers}.
\end{proof}

The transport term of \eqref{eq:proximal-primal} is linear in $\mom$, and both penalties are convex quadratics in $\mom$, so \eqref{eq:proximal-primal} is a weighted projection of a shifted target onto $\momset$.
Its minimizer always exists and is unique in $\mom$; it is a genuine partition, namely the power diagram of the generators $\dual_i = \qq_i + \shift_i$, whenever these generators are distinct.

\begin{proposition}[Derivatives and concavity]
\label{prop:proximal-derivatives}
Wherever the generators $\dual_i$ are distinct,
\begin{subequations}
  \label{eq:proximal-gradient}
  \begin{empheq}[left={\empheqlbrace\,}]{align}
    \partial_{\shift_i} \regdual &= 2\,\cellvol_i (\tcentroid_i - \centroid_i) - \alpha\,\pvol_i\,\shift_i , \\
    \partial_{\weight_i} \regdual &= \pvol_i - \cellvol_i - \beta\,\pvol_i\,\weight_i ,
  \end{empheq}
\end{subequations}
and the Hessian is the facet Laplacian of \eqref{eq:appendix-block} shifted by the diagonal $-\mathrm{diag}(\alpha\pvol_i\Id_d, \beta\pvol_i)$.
Hence $\regdual$ is strongly concave, and its maximizer is unique.
With $\beta = 0$, the only null mode is $\delta\weight_i \equiv c$; with $\alpha = 0$, the null space is that of the Hessian of $\centroiddual$, characterized in Remark~\ref{rem:centroid-kernel}.
\end{proposition}

\begin{proof}[Sketch]
The gradient and Hessian follow from \S\ref{sec:joint-constraints}.
For $\alpha, \beta > 0$ the regularizer is strongly concave and $\jointdual$ is concave.
For $\beta = 0$, a null vector must have $\delta\shift = 0$, and the remaining weight block is the volume Laplacian \eqref{eq:newton-step}, whose kernel is the constants.
For $\alpha = 0$, a null vector must have $\delta\weight = 0$, and the remaining shift block is the Hessian of $\centroiddual$.
\end{proof}

\begin{remark}[Kernel of the Hessian of $\centroiddual$]
\label{rem:centroid-kernel}
The quadratic form of Appendix~\ref{sec:appendix-centroid} vanishes when the affine fields $\delta\varphi_i$ agree on every facet, which does \emph{not} force them to glue into one global affine function.
Agreement on a facet only says that $\delta\varphi_i - \delta\varphi_j$ vanishes on the hyperplane of $\face_{ij}$, \ie that this hyperplane does not move to first order.
The generator perturbations that fix every facet hyperplane are the parallel redrawings of the dual triangulation; for generic connectivity these are only the $(d\!+\!2)$-dimensional redundancy family $\delta\dual_i = s\,\dual_i + \translation$, $\delta\theta_i = s\,\theta_i + c$, and we assume this below.
Within $\centroiddual$, $\delta\theta_i = -2\,\delta\shift_i\tr\tcentroid_i$, so a null direction exists iff the homogeneous linear system $s\,\gamma_i + 2\,\translation\tr\tcentroid_i + c = 0$ for all $i$ has a nonzero solution $(s, \translation, c)$, where $\gamma_i = \theta_i + 2\dual_i\tr\tcentroid_i = \weight_i - \norm{\qq_i - \tcentroid_i}^2 + \norm{\tcentroid_i}^2$ depends only on the data.
Hence the Hessian of $\centroiddual$ is negative definite iff the $n$ vectors $(\gamma_i, \tcentroid_i, 1) \in \R^{d+2}$ span $\R^{d+2}$, which requires $n \ge d + 2$ and is generic, but is strictly stronger than the affine-span assumption of \S\ref{sec:centroid-constraints}.
In the centroidal case $\qq = \tcentroid$, $\weight = 0$, we get $\gamma_i = \norm{\tcentroid_i}^2$, so $\barcentroiddual$ is degenerate exactly when the targets lie on a common sphere or hyperplane (always the case for $n \le d + 1$; \eg also for sites on a regular polygon).
The same count shows that, for $n \ge d + 2$ and non-degenerate data, the set of centroids realizable with free volumes is open (implicit function theorem on $\nabla\centroiddual = 0$).
\end{remark}

The regularizer pulls the generators $\dual_i = \qq_i + \shift_i$ toward the fixed sites $\qq_i$, which are distinct, and so prevents the collapse $\gen \to 0$ of Proposition~\ref{prop:joint-structure}.
Quantitatively, if every target centroid lies in $\domain$, then \eqref{eq:proximal-multipliers} gives $\norm{\shift_i} \le 2\,\domainVol\,\mathrm{diam}(\domain) / (\alpha \min_j \pvol_j)$, so the generators are distinct, and the optimum is a power diagram, as soon as
\begin{equation}
  \alpha > \frac{4\,\domainVol\,\mathrm{diam}(\domain)}{\min_j \pvol_j \, \min_{i \neq j} \norm{\qq_i - \qq_j}} .
  \label{eq:alpha-bound}
\end{equation}
The parameters interpolate between known problems:
$\alpha \to \infty$ drops the centroid term and recovers the volume-constrained power diagram of the sites $\qq$ (\S\ref{sec:background}, with $\beta \to 0$);
$\beta \to \infty$ with $\alpha \to 0$ recovers $\centroiddual$ with zero weights (\S\ref{sec:centroid-constraints});
and $\alpha, \beta \to 0$ recovers the degenerate $\jointdual$.

\paragraph{Volumes hard, centroids soft ($\beta \to 0$).}
With $\cellvol_i = \pvol_i$, the transport term equals $-2\sum_i \pvol_i \qq_i\tr \centroid_i$ up to a constant, and \eqref{eq:proximal-primal} becomes
\begin{equation}
  \min_{\centroid \in \convset(\pvol)} \ \sum_i \pvol_i \norm{\centroid_i - \tgt_i}^2 ,
  \qquad
  \tgt_i = \tcentroid_i + \tfrac{\alpha}{2}\,\qq_i ,
  \label{eq:volume-hard-projection}
\end{equation}
the weighted projection of the shifted target $\tgt$ onto the convex set $\convset(\pvol)$.
The generators are $\dual_i = \tfrac{2}{\alpha}(\tgt_i - \centroid_i)$, \ie the outward normal of $\convset(\pvol)$ at the projection, consistent with \citet[Lemma~4.1]{Bourne2026ConvexOrder}.
If $\tgt$ lies inside $\convset(\pvol)$, the projection is $\tgt$ itself, all generators coincide, and the optimum is a fractional partition, not a power diagram.
The sites $\qq$ are what pushes the target outward: in the centroidal setting $\qq_i = \tcentroid_i$ of $\barjointdual$, the target is the dilation $(1 + \tfrac{\alpha}{2})\tcentroid$, which is exactly the $\lambda$-dilation trick of \citet[\S5.2]{Bourne2026ConvexOrder} with $\lambda = 1 + \tfrac{\alpha}{2}$.
Their Proposition~5.2 then applies: if $\tcentroid$ is realizable by an irreducible power diagram, then for $\alpha$ small enough the maximizer of $\regdual$ (with $\beta \to 0$) is a power diagram with volumes $\pvol$ whose centroids tend to $\tcentroid$ as $\alpha \to 0$.

\emph{Example.}
In $[0,1]$ with $\pvol = (\tfrac12, \tfrac12)$, $\qq = \tcentroid = (0.3, 0.7)$, the projection of $\tgt$ onto $\convset$ is $\centroid_1 = 0.3 - 0.1\alpha$.
For $\alpha < \tfrac12$ it is interior, and the generators coincide ($\dual = (0.5, 0.5)$ at $\alpha = 0.4$).
For $\alpha \ge \tfrac12$ it is the realizable endpoint $\centroid = (\tfrac14, \tfrac34)$ with distinct generators ($\dual = (0.4, 0.6)$ at $\alpha = 1$).
For realizable targets $\tcentroid = (\tfrac14, \tfrac34)$, every $\alpha > 0$ returns the exact diagram.

\paragraph{Centroids hard, volumes soft ($\alpha \to 0$).}
The feasible set of \eqref{eq:proximal-primal} becomes the convex set of admissible volumes $\{\cellvol : ((\cellvol_i\tcentroid_i, \cellvol_i))_i \in \momset\}$, which is contained in the polytope of \S\ref{sec:feasibility-hard}.
The objective is the linear transport term (which favours the relative boundary, \ie power diagrams) plus the volume penalty (which may pull the optimum inside).
By Remark~\ref{rem:centroid-kernel}, realizable centroids with free volumes form an open set for $n \ge d + 2$ (consistent with the count $n(d+1) - d - 2 \ge nd$), so small perturbations of realizable targets stay realizable, but arbitrary targets need not be.
For $n < d + 2$ this fails: in the 1D example above the identities force $\cellvol = (\tfrac12, \tfrac12)$, which together with $\tcentroid = (0.3, 0.7)$ is an interior moment, so no power diagram exists.

\subsection{Exact recovery by proximal point}
\label{sec:proximal-point}

Replacing the fixed prox-center $0$ in \eqref{eq:proximal-dual} by the previous iterate gives the proximal-point method on $\jointdual$:
\begin{equation}
  (\shift^{k+1}, \weight^{k+1})
  = \argmax_{\shift, \weight} \
  \jointdual(\shift, \weight)
  - \frac{\alpha}{2} \sum_i \pvol_i \norm{\shift_i - \shift_i^k}^2
  - \frac{\beta}{2} \sum_i \pvol_i \, (\weight_i - \weight_i^k)^2 .
  \label{eq:proximal-point}
\end{equation}
Writing $\shift = \shift^k + \shift'$ and completing the square in \eqref{eq:shifted-site-weight} shows that each step is \eqref{eq:proximal-dual} with the sites $\qq$ replaced by the current generators $\dual^k$ and the weight regularizer recentred at the correspondingly offset weights; the method is a \emph{re-seeding} loop in which every subproblem is strongly concave.
It is the method of multipliers (augmented Lagrangian) for the jointly constrained primal \eqref{eq:joint-lagrangian}.
By the classical theory of the proximal-point algorithm for concave functions, the iterates converge to a maximizer of $\jointdual$ when one exists (\ie when $\tmom \in \momset$) and diverge otherwise.
A fixed point has $\shift' = 0$ and unchanged weights, so by \eqref{eq:proximal-multipliers} its cells match the prescribed volumes and centroids exactly.
If $\tmom$ lies in $\momset$ but is not realizable (relative interior, or a non-exposed boundary point), the only maximizers are degenerate and the iterates converge to coinciding generators, which is the correct certificate of infeasibility.

\begin{remark}[Conjecture: non-degenerate limit]
By Proposition~\ref{prop:joint-structure}, when $\tmom$ is realizable by a diagram with generators $\dual^\star$ (unique up to the redundancy of \citet[Lemma~2.1]{Bourne2025Inverting}), the maximizers of $\jointdual$ form the half-line $\{ t\gen^\star : t \ge 0 \}$ plus the gauge modes.
The proximal iterates converge to some point of this set, and the limit is non-degenerate if and only if $t > 0$.
We conjecture that $t > 0$ whenever the initial sites are positively correlated with the realizing generators, $\sum_i \pvol_i (\qq_i - \bar\qq)\tr (\dual^\star_i - \bar\dual^\star) > 0$ (bars denoting $\pvol$-weighted means).
For $\qq = \tcentroid$ this holds, as the volume-hard analysis above suggests, since the centroids of a power diagram are positively correlated with its generators by cyclical monotonicity \citep[Lemma~4.10]{Bourne2025Inverting}.
\end{remark}

\subsection{Connection to the Lagrangian flow}
\label{sec:joint-flow}

For fluid particles with fixed volumes $\pvol$, the set of particle positions $\centroid$ admitting \emph{some} power diagram with volumes $\pvol$ and centroids $\centroid$ is the set of exposed points of $\convset(\pvol)$, a dense part of its relative boundary of dimension $nd - d - 1$.
Here the identities \eqref{eq:partition-feasibility} are genuine hard constraints on the unknowns: they fix the center of mass of the fluid.
But the remaining condition is a single scalar, so mere existence of a volume-preserving diagram constrains the particles far less than incompressibility, which needs $n\!-\!1$ independent volume conditions.
This is why the flow needs a rule that ties the cells to the particles (the $\centroiddual$ selection of the flow section), rather than joint realizability.

\subsection{Status}
\label{sec:joint-status}

\emph{Proven (modulo standard regularity):} Proposition~\ref{prop:joint-structure}, Proposition~\ref{prop:proximal-primal}, Proposition~\ref{prop:proximal-derivatives}, the bound \eqref{eq:alpha-bound}, the reduction \eqref{eq:volume-hard-projection}, the equivalence with the dilation trick for $\qq = \tcentroid$ and its small-$\alpha$ guarantee \citep[Prop.~5.2]{Bourne2026ConvexOrder}, and convergence of the proximal point method to some maximizer of $\jointdual$.
\emph{Proven under a genericity assumption, correcting Appendix~\ref{sec:appendix-centroid}:} Remark~\ref{rem:centroid-kernel}, assuming the dual triangulation has no nontrivial parallel redrawings.
\emph{Conjectural:} the non-degenerate-limit condition of the remark, global convergence of a damped Newton method on $\regdual$ (by analogy with the volume-only case), and quantitative bounds on $\alpha$ for exact recovery beyond the $\qq = \tcentroid$ case.

\ifSubfilesClassLoaded{%
  \bibliographystyle{ACM-Reference-Format}
  \bibliography{references}
}{}

\end{document}
